\section{Discussion and Conclusion}
\label{sec:discussion}

\subsection{Orientation is the signal: why grid pooling is load-bearing}

The central deep-learning finding of this work is that reading the Doppler angle
$\theta$ from a B-mode image is an orientation-regression problem, and that the
conventional global-average-pooling (GAP) head is structurally unsuited to it.
GAP collapses a convolutional feature map by averaging over every spatial
position, a reduction that is approximately invariant to in-plane rotation of the
input. Invariance to orientation is the wrong inductive bias when the quantity to
be predicted is itself an orientation, because the head discards the very signal
the target depends on. The effect is large. Under image-level five-fold
cross-validation a frozen DenseNet201 with a GAP head scores
\SI{10.85}{\percent} MAPE, and under the stricter patient-level protocol it scores
\SI{18.7}{\percent}.

Replacing GAP with an orientation-preserving grid-pooling head recovers most of
that loss. The head average-pools the feature map onto a small $G\times G$ grid
and flattens it, so the coarse spatial layout of vessel structure survives into
the regressor. At image level this takes frozen DenseNet201~\citep{Huang2017} from
\SI{10.85}{\percent} to \SI{4.58}{\percent} MAPE under five-fold cross-validation,
roughly halving the error, and reaches \SI{5.84}{\percent} MAPE
(\SI{3.77}{\degree} MAE) on the original single augmented split
(Table~\ref{tab:replication}), reproducing the regime of the EMBC~2019
study~\citep{Patil2019} with no fine-tuning. The backbone, even frozen, already
encodes the vessel orientation; the work is in not pooling it away. The Grad-CAM
attributions~\citep{Selvaraju2017} (Fig.~\ref{fig:gradcam}) are consistent with
this reading, concentrating on the vessel region that defines the flow axis rather
than on background texture. We note that a hand-built structure-tensor prior reaches
\SI{3.16}{\degree} MAE on the narrow base-angle band with no learning at all
(\S\ref{sec:discussion}), so grid pooling is not the only route to a usable angle
estimate; what it provides is a frozen-backbone learned pipeline that is competitive
on the angular target. We offer the pooling lesson as a hypothesis worth testing
beyond this one task: when a regression target is a spatial pose or orientation, the
pooling operator should be chosen to preserve the geometry the label encodes rather
than average it out. The present evidence for that lesson is a single carotid task on
a few dozen subjects.

\subsection{Why newer backbones lose on tiny data}

The architecture bake-off (Fig.~\ref{fig:bakeoff}, frozen, patient-level five-fold
cross-validation, untuned heads) runs against the usual ImageNet-leaderboard
intuition. On the extraction-matched comparison, frozen DenseNet201 leads at
\SI{14.1}{\percent} MAPE ($\pm$\SI{2.19}{\percent}), ahead of the best modern
encoder, ConvNeXt-Base~\citep{Liu2022}, at \SI{15.7}{\percent}
($\pm$\SI{1.75}{\percent}); ConvNeXt-Tiny follows at \SI{16.1}{\percent} and the
EfficientNet and EfficientNetV2~\citep{Tan2019} families trail at
\SIrange{17}{21}{\percent}. The DenseNet and ConvNeXt-Base intervals overlap, so
the ranking among the leaders should be read as a draw rather than a clear win.
The remaining classic backbones (ResNet50~\citep{He2016},
Xception~\citep{Chollet2017}, InceptionV3~\citep{Szegedy2016},
VGG19~\citep{Simonyan2015}) were extracted with the replication-era pipeline, so we
do not rank them against ConvNeXt strictly like-for-like. The claim here is narrow:
on this task no newer or larger backbone improves on a frozen 2017 DenseNet201, not
that classic architectures beat modern ones as a family.

The likely reason is leverage rather than encoder quality. With only \num{84} base
images and frozen weights, no gradient reaches the backbone. The encoder supplies a
fixed feature basis, and the only fitted parameters live in the shallow head.
Architectures whose ImageNet advantage is realized through end-to-end training have
nothing to grip under a frozen protocol, and their feature statistics, tuned to a
classification objective, are no better matched to an angular regression target than
those of an older network. DenseNet201 happens to fit this task well across our
measurements: it yields the highest $R^2$, wins the replication, and remains
strongest after tuning. We read this as an empirical fact about our extraction
rather than a proven property of dense feature reuse, and we note that the original
study found VGG19 best under per-model tuning. We carry DenseNet201 forward. The
caution for small-cohort imaging is that benchmark rankings established on
million-image corpora do not transfer to a frozen-feature regime over a few dozen
subjects.

\subsection{Two protocols, two questions}

We report and tune the estimator under two complementary sampling protocols, which
answer different questions. \emph{Image-level sampling} partitions the
rotation-augmented corpus at the image level, the original study's protocol, and
measures how accurately the model reads $\theta$ across the population of
orientations and imaging conditions the augmentation spans. \emph{Patient-level
sampling} holds out whole volunteers via grouped cross-validation and measures
generalization to anatomy the model has never seen. The patient-level number is
naturally larger (Table~\ref{tab:dual_protocol}, Fig.~\ref{fig:protocol_gap}): once
volunteer-specific appearance can no longer be exploited, the model must transfer
across subjects. For the tuned stacked ensemble the image-level error is
\SI{2.79}{\percent} MAPE (\SI{1.96}{\degree} MAE, $R^2=0.995$) and the
patient-level error is \SI{8.53}{\percent} MAPE (\SI{5.93}{\degree} MAE,
$R^2=0.952$). The patient-level figure carries appreciable fold variance: per-fold
MAPE standard deviations are \SIrange{2.6}{4.4}{\percent} on the few held-out groups
per fold (Table~\ref{tab:dual_protocol}), so a spread of a few percentage points
around the \SI{8.53}{\percent} headline is expected, and the cross-subject estimate
should be read with that uncertainty in mind.

The gap between the two protocols is itself a measurement. It quantifies how much of
the within-population accuracy is anatomy-specific, that is, the cost of
cross-subject generalization. A reader sizing the method for deployment should weigh
the patient-level figure when the operative question is a new patient, and the
image-level figure when it is accuracy across the imaging conditions the corpus
represents. The per-backbone ordering is broadly stable across both lenses, though
not perfectly: under patient-level MAPE, VGG19 (\SI{14.97}{\percent}) and
InceptionV3 (\SI{14.66}{\percent}) sit within \SI{0.3}{} percentage points and
exchange ranks. The lesson for small-cohort medical imaging is that the sampling
protocol is part of the result. Augmentation manufactures correlated samples, so the
unit of statistical independence is the patient, not the image, and reporting both
protocols is more informative than privileging one.

Two comparisons need a word on aggregation, because the numbers are aggregated
differently. The single tuned DenseNet201 (\SI{3.00}{\degree} MAE,
\SI{4.03}{\percent} MAPE image-level) reproduces the best single model of the
EMBC~2019 work~\citep{Patil2019} ($\approx\SI{2.87}{\degree}$ MAE,
\SI{4.03}{\percent} MAPE). Head tuning~\citep{Akiba2019,Bergstra2011} and a
five-model stacked ensemble~\citep{Wolpert1992} improve on it and additionally
supply the previously unmeasured patient-level result. The per-backbone member
figures above are per-fold cross-validation means, whereas the ensemble figures are
pooled out-of-fold. Aggregated the same pooled-OOF way, the single tuned
DenseNet201 reaches \SI{7.80}{\degree} MAE patient-level, so the genuine ensembling
gain over the best single model is about \SI{1.9}{\degree} MAE; the larger apparent
gap between member and ensemble rows partly reflects the per-fold-mean versus
pooled-OOF aggregation rather than the ensemble alone.

\subsection{Uncertainty quantification and method-vs-reference agreement}

A point estimate alone is not clinically actionable, so we also report calibrated
prediction intervals. Split-conformal intervals~\citep{Vovk2005,Lei2018}, computed
on a single patient-disjoint calibration/test split, attain empirical coverage at or
above their nominal level at every level tested (Table~\ref{tab:clinical},
Fig.~\ref{fig:calibration}): the \SI{90}{\percent} interval has half-width
$\pm\SI{20.50}{\degree}$ with \SI{95.2}{\percent} empirical coverage, and the
\SI{95}{\percent} interval has half-width $\pm\SI{26.03}{\degree}$ with
\SI{97.8}{\percent} coverage. Two qualifications apply. The validity guarantee is
distribution-free and finite-sample under exchangeability; with grouped data and
only ten volunteers, exchangeability across patients is the assumption most at risk,
and coverage is marginal over the population rather than conditional on the
individual case. The reported coverages are single empirical numbers on a small test
set and therefore carry wide binomial uncertainty, so the observed over-coverage is
as consistent with conservative (somewhat too wide) intervals as with tight
calibration.

The intervals are wide, and the width is a property of a ten-volunteer cohort rather
than a modelling artifact. A wide interval that covers is still more useful for a
velocity-correction decision than a narrow one that does not, because it does not
mask the state of the evidence. We are careful about what such an interval offers in
practice: because the half-width is fixed across cases, it reports population-level
coverage and does not flag which individual readings are unreliable. For the
$\cos\theta$ velocity correction a half-width of \SIrange{20.50}{26.03}{\degree} can
move the correction factor markedly, so whether a band this wide is directly
actionable per case is an open question that external data would have to settle.

The Bland--Altman analysis~\citep{BlandAltman1986} (Fig.~\ref{fig:bland_altman})
reports a bias of $-\SI{4.31}{\degree}$ (method minus reference: the model reads
slightly low) and limits of agreement of $-24.25$ to $+\SI{15.63}{\degree}$. There
is exactly one human angle reading per image, the MATLAB-GUI annotation that serves
as the label, so this is a method-vs-reference comparison against a single reading
and not an inter-observer study. Without a second reader we cannot quantify the
irreducible human labelling variability against which an automated method should
ultimately be judged, and the reference reading is itself a noisy estimate of the
true flow angle. The bias also has a direction. A consistent low-angle bias shifts
the $\cos\theta$ velocity correction the same way for every case, which for an
application like stenosis grading is a systematic error rather than random scatter,
and a directional offset of this size warrants attention even though it is small and
fully characterized in our cohort.

Two further analyses are encouraging within these bounds, though on a different
evaluation unit. Both are computed on the \num{84} base images over the narrow
base-angle band, not on the patient-level OOF protocol, so they are not directly
comparable to the patient-level ensemble MAE of \SI{5.93}{\degree}. Rotation-based
test-time augmentation, combined under a seam-safe circular median over de-rotated
views, reduces the raw base-image MAE from \SI{7.80}{\degree} to \SI{4.72}{\degree}
with no retraining. A circular fusion of the learned estimate with a classical
structure-tensor angle prior reaches \SI{2.72}{\degree} MAE on the same band,
beating either the learned model or the geometric prior (\SI{3.16}{\degree} MAE)
alone, which suggests that a data-driven encoder and a hand-crafted geometric cue
capture partly complementary information.

\subsection{Clinical relevance and deployment}

The clinical appeal of an image-only estimator is that it slots into the existing
Doppler workflow without changing it. Angle correction is today a manual cursor
placement, a recurring source of error and, as noted in \S\ref{sec:intro}, a
deficiency flagged in vascular-laboratory accreditation. An automated reading
accurate to a few degrees could pre-populate or sanity-check that cursor and reduce
the operator-dependent variation in the measured angle. We are careful not to
overstate this. The labels come from a single annotator with no inter-observer
baseline, so an estimator trained against them can at best reproduce that
annotator's readings and their biases; it cannot independently correct operator
error or claim to standardize against a validated ground truth. Any downstream
benefit to velocity estimation or stenosis grading is a plausible motivation rather
than a measured result, since neither was evaluated here.

Because the estimator consumes only the formed grayscale B-mode frame, with no
color Doppler, no segmentation, and no knowledge of the vendor-specific beamforming
path, it could in principle be packaged as a software add-on. We have not tested
this: the cohort is single-scanner and single-center, so cross-vendor behavior is
unknown and the device-agnostic framing is a design property rather than a validated
capability. The conformal intervals are relevant to deployment in that a tool
reporting a calibrated band alongside its point estimate communicates population-level
uncertainty rather than a bare number, which is a safer default than silent
over-confidence. As noted above, a fixed-width marginal interval does not identify
which individual cases are uncertain, so it supports a general posture of caution
more than per-case triage. Realizing any of this would require the external
validation discussed below; the present results establish that the core estimate is
accurate and reproducible under our protocols, which is what justifies pursuing that
validation.

\subsection{Limitations}

The cohort is small: about ten volunteers and \num{84} base images from a single
center, a single scanner, and the carotid only. Patient-level cross-validation
therefore rests on a handful of held-out subjects per fold, the per-fold variance is
non-negligible (see the \SIrange{2.6}{4.4}{\percent} MAPE standard deviations of
Table~\ref{tab:dual_protocol}), and external validity to other vendors, anatomies,
and acquisition settings is untested. The method has never seen a tortuous vessel, a
bifurcation, plaque, or a non-carotid site, so any workflow claim beyond the imaged
population is speculative. There is a single reference reading per image, so the
agreement analysis is method-vs-reference rather than inter-observer, and the upper
bound on achievable accuracy set by human labelling noise is unknown. Finally, every
learned result here uses a frozen ImageNet backbone with only a trained pooling head;
the encoder is never updated. The estimator should be read as a strong, fully
reproducible frozen-feature baseline, not as the ceiling of what this task admits.

\subsection{Future work}

The frozen-only scope leaves several directions for once dedicated accelerator
hardware is available. End-to-end fine-tuning of the backbone, left unrun here, is
the most direct candidate for improving the patient-level number, since it would let
the encoder adapt its features from a generic classification basis to the
angular-regression objective. Self-supervised pretraining on the unlabelled carotid
frames could yield representations less dependent on the small labelled cohort and so
better cross-subject transfer. Geometry-aware alternatives to grid pooling, an
explicitly circular treatment of the angular target across the full $360^{\circ}$
range, and orientation-equivariant heads are natural architectural extensions of the
pooling insight. On the evaluation side, a multi-reader study would convert the
method-vs-reference agreement into a proper inter-observer comparison and establish
the human-noise floor, and a multi-center, multi-vendor cohort would test the
external validity that a single-center dataset cannot. None of these is claimed as a
result of the present paper; they are the experiments the present results motivate.

\subsection{Conclusion}

We have presented a reproducible, frozen-backbone estimator of the Doppler
beam-to-vessel angle from a single grayscale B-mode carotid image, and isolated the
design choice that makes a frozen pipeline competitive on this orientation-regression
task: orientation-preserving grid pooling in place of global average pooling. Optuna
head tuning and a five-model stacked ensemble yield our best estimator, reported
under two complementary sampling protocols. The image-level error is
\SI{2.79}{\percent} MAPE (\SI{1.96}{\degree} MAE, $R^2=0.995$) and the patient-level
error is \SI{8.53}{\percent} MAPE (\SI{5.93}{\degree} MAE, $R^2=0.952$), the latter
with the fold variance noted above. Against the one prior point of comparison on this
task, the EMBC~2019 study, the tuned ensemble improves on the published image-level
figures. An architecture bake-off shows that no newer or larger backbone improves on
a frozen 2017 DenseNet201 on this cohort, and the evaluation equips the method with
distribution-free conformal intervals and a method-vs-reference agreement analysis.
The takeaway is twofold. For this task, automated Doppler-angle estimation is
accurate and now reproducible, characterized under both the image-level and the
stricter patient-level lens. For small-cohort medical imaging more broadly, matching
the pooling operator to the geometry of the target, reporting accuracy under more
than one sampling protocol, and supplying calibrated intervals rather than a single
point estimate are practices that make such a method easier to trust.

\subsection*{Reproducibility}

Every table and figure in this paper is regenerated from saved results by code, with
no hand-entered numbers. The data is the public SPLab common-carotid B-mode set; the
pipeline (data preparation, the rotation-augmentation corpus, frozen feature
extraction, the grid-pooling head, Optuna tuning, ensembling, and the conformal /
Bland--Altman / TTA evaluation) is deterministic under a fixed seed ($42$) and runs
on a CPU. After preparing the corpus and cross-validation results, the analysis
tables are produced by a single table-generation script and the figures by a single
figure-generation script; both read only the saved cross-validation summaries and
out-of-fold predictions. The frozen patient- and image-level five-fold
cross-validation that underlies the protocol comparison and the architecture bake-off
is likewise produced by a single command per backbone, so the reported means and
standard deviations can be reproduced end to end from the raw images.
